\chapter{Perspectives on Caste: G. S. Ghurye}

\section{The Four Scholars and the Beginnings of Caste Writing}

\begin{KeyConcept}
\begin{itemize}
\item Section B covers \textbf{Perspectives on Caste}, built around \textbf{four key scholars}: \textbf{G. S. Ghurye}, \textbf{M. N. Srinivas}, \textbf{Andre Beteille} and \textbf{Louis Dumont}.
\item The first writings on caste can be traced to the \textbf{1820s}. One of the first books on caste (with origin, evolution and the modern scenario) is attributed to \textbf{James Stuart Mill (1824)} --- and \textbf{Ghurye copied this structure}: origin $\rightarrow$ evolution $\rightarrow$ present scenario $\rightarrow$ the future (will India become \textbf{caste-less}, as the makers of the Constitution envisioned, or \textbf{plural} --- with more and more castes?) $\rightarrow$ and \textbf{caste outside India}.
\end{itemize}
\end{KeyConcept}

\begin{ExampleBox}
\begin{itemize}
\item \textbf{Satish Deshpande's sociological reading of a joke} (from one of his memorial lectures): when OBC reservation was first implemented (1991--92), a joke went viral --- of 10 astronauts sent to explore Mars/Moon, 6 were OBC, 2 SC, 1 ST, and 1 general (if space was left).
\item The implication: for the \textbf{lower caste}, \textbf{caste overrides qualification}; for the \textbf{upper caste}, \textbf{qualification overrides caste}. Caste is \textbf{hyper-visible among the lower castes and invisible among the upper castes} --- upper castes are spoken of as 'general', as if they are caste-less, although they have benefited most from caste.
\end{itemize}
\end{ExampleBox}

\section{The Four Sources of Ghurye's Ideas on Caste}

\begin{KeyConcept}
\begin{itemize}
\item Ghurye's ideas on caste can be traced to \textbf{four sources/factors}:
\item 1. \textbf{Indological sources} --- the Vedic texts, Samhitas and Brahmanas, i.e. how caste was described in the ancient Sanskrit classical texts (Ghurye works like an archaeologist of history, uncovering when caste arose and how it was practised).
\item 2. \textbf{The Indo-Aryan Racial Theory} (of \textbf{Risley}) --- caste developed from race.
\item 3. \textbf{The idea of the sub-caste} --- a sub-caste is itself a caste (it is 'sub' only in relation to a caste above it, and it has sub-sub-castes below it; no caste is a single homogeneous unit --- it has internal factions/segments).
\item 4. \textbf{Caste outside India}.
\item From these four sources come Ghurye's \textbf{six features of caste}.
\end{itemize}
\end{KeyConcept}

\section{The Indo-Aryan Racial Theory}

\begin{KeyConcept}
\begin{itemize}
\item \textbf{Indo-Aryan Racial Theory} (given by \textbf{Risley}): caste developed from \textbf{race}, and \textbf{occupation stabilised it}.
\item The \textbf{Aryans} (who came from outside, from the Persian side) were \textbf{fair}; the \textbf{non-Aryans} (the indigenous people, called \textbf{Dasas}) were \textbf{dark}. Fair and dark = \textbf{varna} (Sanskrit for 'colour').
\item To preserve their unique identity and purity, the Aryans imposed \textbf{restrictions on intermarriage} and on \textbf{occupational purity}; the groups they encountered were excluded from these occupations.
\item The Aryans already practised a three-class system (Brahmin, Kshatriya, Vaishya --- by other names) in Persia; in India they continued it, with a \textbf{fourth class outside} them. Each class had a specific occupation. Restricting intermarriage with the outsiders automatically produced a social group marked by \textbf{endogamy and hierarchy --- the caste}.
\end{itemize}
\end{KeyConcept}

\begin{CriticalPoint}
\begin{itemize}
\item \textbf{Ghurye's own verdict} (a line to remember verbatim): the racial theory is \textbf{'valid only for Hindustan proper'} --- i.e. only in \textbf{North India} (the Hindo-Gangetic plain). In \textbf{South India, endogamy was not practised} and there was no restriction on intermarriage. So the theory is \textbf{partially valid, partially invalid}.
\item The race--caste correlation \textbf{cannot be scientifically verified}: geographical factors are ignored --- in Kashmir even Dalits are fair, and in Kanyakumari even Brahmins are dark.
\item To verify it, Ghurye conducted \textbf{anthropometric analysis} (measurements of the \textbf{nasal index} --- the kind of nose/brow measurements used in crime documentaries), an ideal-type model inspired by \textbf{Lombroso} (of 'Born Criminal' fame) and \textbf{Risley} --- the first to do such measurements. (This detail is not exam-important.)
\end{itemize}
\end{CriticalPoint}

\section{The Indological Account: Four Phases of Caste}

These points are in Ghurye's book \textbf{'Caste and Race' (1932)}: chapters 1--2 = Indo-Aryan racial theory; chapters 3--4 = the Indological account of caste, written as \textbf{four temporal phases of caste}.

\begin{KeyConcept}
\begin{itemize}
\item \textbf{Phase 1 --- Vedic period (Early Vedic; 1500 BC -- 920 BC)}: society was \textbf{egalitarian and tribal}; the order was the \textbf{varna vyavastha} (not caste/jati), where \textbf{anyone could change his varna}. The \textbf{Purusha Sukta of the Rig Veda} mentions only \textbf{three classes} (Brahmin, Kshatriya, Vaishya/Vis) --- no fourth, and no restriction; there was \textbf{no Brahminical supremacy} and no question of higher/lower or untouchable.
\item \textbf{Phase 2 --- Post-Vedic period (920 BC -- 400 BC)}: dominated by the sacred laws of the Aryas, the great epics and Gita literature. The \textbf{four varnas were firmly established} --- the \textbf{fourth varna} emerged from the \textbf{intermixing of the three Aryan classes with the non-Aryans} (the outsiders). It was seen as the most impure; those born of intermixing between the third and the non-Aryan were the most polluted. Mixed castes are called \textbf{'Sankar Jati'} in Hindu literature (Kaliyuga is called the phase of Sankar Jati --- when caste order breaks down). When it was realised that intermixing weakens caste, \textbf{endogamy came to be strictly followed}, and \textbf{Brahmins gained supremacy}.
\item \textbf{Phase 3 --- the period of the Dharma Shastras / Hindu emperors (200 BC onwards)}: the \textbf{Dharma Shastras} (written sacred texts, like a law code governing everyday life) were \textbf{codified} --- for the first time in writing: who may marry whom, who may eat/sit with whom. Rules on \textbf{social interaction}, restrictions on \textbf{inter-dining and inter-marriage}, and --- most important --- \textbf{different punishment for the same crime committed by different castes} (differential treatment). Summed up in the \textbf{Manu Smriti}. The Brahmin was projected as the lord of creation, to be given gifts (dana). Caste became \textbf{completely rigid}.
\item \textbf{Phase 4 --- the modern period (18th century onwards, when the British came)}: caste was \textbf{neither flexible nor rigid, but both}.
\end{itemize}
\end{KeyConcept}

\begin{KeyConcept}
\begin{itemize}
\item \textbf{How caste became flexible} in the modern period: the \textbf{social-religious reform movements} attacked caste, and \textbf{British legislation} --- the \textbf{Caste Disabilities Removal Act 1850}, the \textbf{Widow Remarriage Act 1856}, and the \textbf{Special Marriage Act 1872} --- weakened the caste system.
\item \textbf{How caste became rigid}: \textbf{reservation} (caste-based reservation makes more castes demand reservation --- people 'become more caste'); and the \textbf{census} (classifying people by caste reproduces caste consciousness rather than eroding it). Both strengthened caste.
\item The overall net effect (weakening vs strengthening) is \textbf{not measurable} --- both happened, for different reasons.
\end{itemize}
\end{KeyConcept}

\begin{KeyConcept}
In conclusion, Ghurye explains the \textbf{evolution of caste} using three references: the \textbf{Indological account}, \textbf{modern historical texts}, and the \textbf{British census}. This is both an \textbf{evolutionary} and a \textbf{comparative} analysis (how caste was vs how caste is).
\end{KeyConcept}

\section{The Book 'Caste and Race': Fission/Fusion, Kin and Politics}

\begin{KeyConcept}
\begin{itemize}
\item 'Caste and Race' (1932) was revised many times; chapters 6--11 were added later, giving the book a new edition structure (later editions titled e.g. 'Caste and Class', 'Caste and Occupation'). The later chapters:
\item \textbf{Chapter 8 --- Fission and Fusion}: \textbf{fission} = the breakdown of a caste into multiple \textbf{sub-castes}; \textbf{fusion} = those sub-castes coming together in a \textbf{coalition} (a caste association) during politics. Ghurye's point: \textbf{caste is not a homogeneous group} --- within it are multiple divisions; even a 'homogeneous' caste becomes heterogeneous sub-castes (endogamy operates at the sub-caste level, not the caste level).
\item \textbf{Chapter 9 --- Caste, sub-caste and kin}: kin = \textbf{kindred} (maternal and paternal relatives --- mama, mami, fufa, fufi, bua, bhanja, bhatiji, cousins). Caste is linked to kinship --- e.g. the \textbf{biradari} (brotherhood/fraternity) in North India; in single-caste villages, caste reproduces generation after generation.
\item \textbf{Chapter 10 --- Caste and politics (general)}.
\item \textbf{Chapter 11 --- Caste and politics in Tamil Nadu} --- showing Ghurye was an \textbf{empiricist} (speaking on data, not speculation). This and the general chapter influenced M. N. Srinivas (dominant caste), Rajni Kothari, and Andre Beteille (who did field study in Tamil Nadu).
\item \textbf{Chapter 12 --- The future of caste in India: caste-less or plural?}
\end{itemize}
\end{KeyConcept}

\section{Caste and Politics: The Census as Creator of Caste}

\begin{KeyConcept}
\begin{itemize}
\item In 'Caste and Politics' (chapter 10), Ghurye examines the relation between caste and politics against the then-contemporary situation. His argument: the \textbf{census is not a passive instrument} merely collecting data --- it has \textbf{objectifying power}: \textbf{in the process of collecting data, it creates new groups} (it creates caste).
\item His deeper claim: \textbf{Indians never practised caste as such --- it was the British who divided them into caste groups through the census.}
\item Example: \textbf{L. Middleton}, the \textbf{census officer of Punjab in 1921}, found that people \textbf{refused to give their caste} --- they said 'we do not have caste' --- yet he imposed a caste on them.
\item Why the British conducted the caste census: they wanted \textbf{data to frame affirmative policies} and give \textbf{adequate representation} (percentage-wise) to communities demanding representation --- i.e. they needed to know each caste's population.
\item Ghurye's verdict: \textbf{the British were trying to cure the problem they themselves had created in the first place.} The census \textbf{manufactured caste and made it an all-India phenomenon} (Ghurye is, like an 'Indian Durkheim', concerned with integration).
\end{itemize}
\end{KeyConcept}

\begin{CriticalPoint}
\begin{itemize}
\item Ghurye was \textbf{against the caste census} for these reasons:
\item 1. It \textbf{divides India} --- it brings out the very caste identity that the Constitution tries to do away with.
\item 2. It \textbf{creates caste} --- every 10 years (the census cycle) the old caste discourse is revived in newspapers and on television.
\item 3. It \textbf{unintentionally creates caste associations} --- e.g. \textbf{Bahujan} = SC + ST + OBC clubbed together into a caste association that became a political party (the \textbf{BSP}).
\item 4. It pushes \textbf{tribals to claim caste identity} --- the \textbf{Hinduization of tribes}: realising Hindus are privileged and tribals isolated, a tribal emulates Hindu customs and becomes a 'caste Hindu' through \textbf{Sanskritization} --- which is a \textbf{gateway/entry into the caste structure} (a process by which the caste-less gain caste).
\item 5. It makes \textbf{Sikhs worry about undercounting} --- the caste census is largely Hindu-specific (Sikhism, like Buddhism under Guru Nanak, was against caste and Brahmanism, but conversions from Hinduism --- mostly by the oppressed sections seeking refuge --- brought caste into Sikhism); they feared not being counted as Indians at all.
\item 6. It leads to the rise of the \textbf{anti-Brahmanical movement} (which began in South India with \textbf{Dravid consciousness}) --- dividing the very India that Ghurye saw as united by a common culture diffused all over India.
\end{itemize}
\end{CriticalPoint}

\begin{KeyConcept}
\begin{itemize}
\item \textbf{Bernard Cohn} and others have made the same point about the census's power to create caste.
\item On opposing the caste census \emph{today}: the reasons differ depending on \textbf{who} is opposing (a university professor vs someone who does not want lower castes to prosper) --- it is a matter of perspective, not a single structural reason.
\end{itemize}
\end{KeyConcept}

\section{Substantialization and 'Caste Patriotism'}

\begin{KeyConcept}
\begin{itemize}
\item Forming associations on the basis of caste \textbf{substantializes caste} --- it makes caste more and more crystallised. This leads to what Ghurye \textbf{coined} as \textbf{'caste patriotism'}.
\item In the anti-Brahman movement, non-Brahmins still have caste (they are against Brahmins, but have their own caste); they form \textbf{coalitions around their own caste identities against the Brahmin identity} --- this is \textbf{caste patriotism}, not patriotism at the Indian level (e.g. the \textbf{Jat agitation/protest} is caste patriotism --- showing solidarity for a common caste cause even if reservation is not obtained).
\item \textbf{Louis Dumont's theory of the 'substantialization of caste'} was actually used \textbf{first by Ghurye}.
\end{itemize}
\end{KeyConcept}

\begin{KeyConcept}
In conclusion, Ghurye is highly concerned about \textbf{fissiparous tendencies} (tendencies that break India apart) --- tendencies that threaten the Hindu culture that (he believes) had diffused all over India and united it. He blames the \textbf{British} (and, in passing, the \textbf{Muslims}) for these fissiparous tendencies.
\end{KeyConcept}

\section{The Six Features of Caste (the framework)}

\begin{KeyConcept}
\begin{itemize}
\item Ghurye's \textbf{last major contribution on caste} is his statement of the \textbf{six features of caste}.
\item Crucially, these are the features of \textbf{caste as it was present in the 1920s} --- i.e. of \textbf{traditional Indian society} (before 1947, the 'British India' that brought modernization). Some may still be valid today, but they describe the 1920s.
\item In other words, Ghurye has given a \textbf{comparative, evolutionary, historical and Indological account of caste}.
\item (The detailed enumeration of the six features is taken up in the next class.)
\end{itemize}
\end{KeyConcept}

\begin{QuickRevision}
\begin{itemize}
\item Four scholars: Ghurye, Srinivas, Beteille, Dumont. First caste book: James Stuart Mill (1824) --- origin$\rightarrow$evolution$\rightarrow$present$\rightarrow$future (caste-less vs plural)$\rightarrow$caste outside India.
\item Ghurye's four sources: Indological texts; Indo-Aryan Racial Theory (Risley); sub-caste; caste outside India $\rightarrow$ six features.
\item Racial theory: Aryan (fair) vs Dasa (dark); caste from race, occupation stabilised it; Ghurye: 'valid only for Hindustan proper' (North); cannot be verified (Kashmir Dalit fair, Kanyakumari Brahmin dark); anthropometric/nasal index (Lombroso/Risley).
\item Four phases (Caste and Race 1932): Vedic (egalitarian, varna vyavastha, Purusha Sukta 3 classes, no Brahmin supremacy); Post-Vedic (4th varna from intermixing, Sankar Jati, endogamy, Brahmin supremacy); Dharma Shastras (codified, Manu Smriti, inter-dining/marriage bans, differential punishment, rigid); Modern (both flexible --- reform movements + British laws 1850/1856/1872 --- and rigid --- reservation + census).
\item Later chapters: Fission/Fusion (sub-castes); Caste, sub-caste, kin (biradari); Caste and politics (general + Tamil Nadu, empiricist) $\rightarrow$ influenced Srinivas/Kothari/Beteille; Future (caste-less or plural).
\item Census: objectifying power --- creates caste (L. Middleton, Punjab 1921); British curing a self-created problem; Ghurye against caste census (divides India, creates caste, caste associations/Bahujan/BSP, Hinduization of tribes via Sanskritization, Sikh undercounting fears, anti-Brahman movement).
\item Substantialization $\rightarrow$ 'caste patriotism' (Ghurye coined; Dumont later); fissiparous tendencies blamed on British (and Muslims).
\item Six features of caste: valid as of the 1920s (traditional society); comparative/evolutionary/historical/Indological account.
\end{itemize}
\end{QuickRevision}

